% Generated draft: requires mathematical and source review.
\documentclass[11pt]{article}
\usepackage[margin=25mm]{geometry}
\usepackage{fontspec}
\setmainfont{Latin Modern Roman}
\usepackage{amsmath,amssymb,mathtools,unicode-math}
\setmathfont{Latin Modern Math}
\usepackage{enumitem,needspace,xurl,hyperref}
\setlength{\parindent}{0pt}
\setlength{\parskip}{5pt}
\setlength{\emergencystretch}{3em}
\newcommand{\N}{\mathbb{N}}
\newcommand{\Z}{\mathbb{Z}}
\newcommand{\Q}{\mathbb{Q}}
\newcommand{\R}{\mathbb{R}}
\newcommand{\C}{\mathbb{C}}
\newcommand{\F}{\mathbb{F}}
\newcommand{\A}{\mathbb{A}}
\newcommand{\PP}{\mathbb{P}}
\newcommand{\E}{\mathbb{E}}
\newcommand{\eps}{\varepsilon}
\newcommand{\dd}{\,\mathrm{d}}
\begin{document}
\section*{UnsolvedMath ID: 11000092}
\textbf{Source problem code:} catalogue entry
\subsection{Definitions and mathematical statement}
In the standard Torelli-space interpretation, \(\mathcal T_g\) parametrizes smooth genus-\(g\) curves with a symplectic homology marking, \(\mathcal T_g^c\) allows stable curves of compact type, and \(\mathcal T_g^{c,\mathrm{red}}\) is its reducible locus. Compact type means that the dual graph is a tree, so the generalized Jacobian is compact. These conventions and the coefficient ring must be confirmed from the chapter; they are not defined in the extracted problem. The source's \(\operatorname{Sp}_g(\Z)\) acts by changing markings (customarily \(\operatorname{Sp}(2g,\Z)\)).

The principal task is to compute the module
\[H_c^k(\mathcal T_g^{c,\mathrm{red}})\]
under this action in a range \(k\geq k_0\), including both the compactly supported cohomology groups and their representation structure. The range threshold \(k_0\), genus range and coefficients are not supplied.

The following reported information is preserved with extraction caveats. The text writes
\[H_c^{6g-7}(\mathcal T_g^{c,\mathrm{red}})=H_0(B_g^c),\qquad
H_c^{6g-8}(\mathcal T_g^{c,\mathrm{red}})\twoheadrightarrow H_1(B_g^c),\]
where \(B_g^c\) is called the normalization, or disjoint union of irreducible components. It describes those components through symplectic translates of the gluing products
\[\mathcal T_{h,1}^c\times\mathcal T_{g-h,1}^c,
\qquad1\leq h\leq\lfloor g/2\rfloor.\]
The extracted orbit/coset indexing is garbled, including possible interchange stabilizers when \(h=g-h\), so it is not replaced by an invented quotient formula. The two displayed degrees also require checking: they appear shifted relative to the later diagram and the usual real dimension \(6g-8\) of this reducible locus. They are source reports, not validated identities or silently corrected statements.

The background folk conjecture says \(H_k(\mathcal T_g)\) is finitely generated for \(k<g-1\), imposing finiteness constraints on \(\mathcal T_g^c\) and its reducible locus if true. The coefficient and degree conventions are not repeated. For \(g\geq3\) the source gives a Gysin sequence whose arrows/line breaks are damaged. Legible fragments include
\[H_c^{6g-10}(\mathcal T_g^{c,\mathrm{red}})\to H_3(\mathcal T_g)\to H_3(\mathcal T_g^c),\]
\[H_c^{6g-9}(\mathcal T_g^{c,\mathrm{red}})\to H_2(\mathcal T_g)\to H_2(\mathcal T_g^c),\]
\[H_c^{6g-8}(\mathcal T_g^{c,\mathrm{red}})\to H_1(\mathcal T_g)
\xrightarrow{\tau}H_1(\mathcal T_g^c)\to0.\]
The diagram additionally displays \(H_c^{6g-10}(\mathcal T_g^c)\), \(H_c^{6g-9}(\mathcal T_g^c)\), \(H_1(B_g^c)\), \(H_0(B_g^c)\) and \((\bigwedge^3 H)/H\), but their full connecting/identification arrows cannot be reconstructed reliably. It identifies \(\tau\) as the Johnson homomorphism induced by \(\mathcal T_g\hookrightarrow\mathcal T_g^c\), with \(H\) the usual first-homology module. No complete exact sequence is fabricated from the fragments.

A further suggested topic is the topology of the branching locus of the period map \(\mathcal T_g^c\to\mathcal J_g^c\), described as the compact-type hyperelliptic locus \(\mathcal H_g^c\). The source invokes A'Campo [1]: the image of the hyperelliptic mapping class group \(\Delta_g\) contains the level-two symplectic subgroup, customarily
\[\operatorname{Sp}_g(\Z)[2]=\{A\in\operatorname{Sp}_g(\Z):A\equiv I\pmod2\}.\]
Its extracted definition repeats the level subgroup inside its own defining set and needs checking. It defines \(\Delta_g\) as the centralizer of a hyperelliptic involution in \(\Gamma_g\), and as the orbifold fundamental group of smooth hyperelliptic moduli. The mod-two image is printed as \(S_{2g-2}\), described as permuting Weierstrass points, while the component formula uses \(S_{2g+2}\):
\[\#\pi_0(\mathcal H_g^c)=\frac{|\operatorname{Sp}_g(\F_2)|}{|S_{2g+2}|}
=\frac{2^{g^2}\prod_{j=1}^{g}(2^{2j}-1)}{(2g+2)!}.\]
The two symmetric-group indices are incompatible as written and require source confirmation; the formula is retained as reported. The source further states that each component is smooth and immerses in Siegel space \(\mathfrak h_g\) by the period map, that the irreducible components of \(\mathcal H_g\) are disjoint in \(\mathcal J_g\), and that in genus three there are \(36\) compact-type components whose images are cut out by the \(36\) even theta nulls \(\vartheta_\alpha:\mathfrak h_3\to\C\). Definitions and genus hypotheses for these reports are missing; they are contextual results, not an independent verification or the principal module computation.
\subsection{Short English statement}
Compute high-degree compactly supported cohomology of the reducible compact-type Torelli locus as a symplectic module. The record also includes finiteness and hyperelliptic-locus context, but its degree shifts, coset indexing and displayed Gysin diagram are damaged.
\subsection{Sources}
\begin{itemize}
\item Local UnsolvedMath registry, permanent ID \(11000092\), exact source code \texttt{catalogue entry}.
\item Farb, \emph{Problems on Mapping Class Groups and Related Topics}, Problem 4.7, PDF page 77; the supplied running text names R. Hain: \url{https://www.math.uchicago.edu/~farb/papers/mcgbook.pdf}. Registry access 2026-07-29; pdf-font-structure extraction. Repeated local numbering is distinguished by page.
\item A'Campo [1] is cited for the hyperelliptic monodromy image; the full entry is absent. The summary mentions unspecified Hain/Torelli-cohomology progress without an exact resolving reference.
\end{itemize}
\end{document}
